% !TeX root = ../教材题库_全部条目_转写工作稿.tex
% Source: XJ-B1-C04-EX-P122-Y1421
% 第4章 / 4.3 对数函数 / 教材 p.122、123 / PDF 第129、130页
% OCR未校对：公式与文字均可能有误；未完成数学TeX转写；禁止直接发给学生
% 原题及相邻原解截图（核对用）：../题目原图/XJ-B1-C04-EX-P122-Y1421-01.webp, ../题目原图/XJ-B1-C04-EX-P122-Y1421-02.webp
% 本文件只在转写完成后作为题干使用；OCR 条目尚不是可靠题干。
\begin{bankitem}{4.3 例题 例题11}{第4章 / 4.3 对数函数 / 教材 p.122、123 / PDF 第129、130页}{XJ-B1-C04-EX-P122-Y1421}{OCR未校对：公式与文字均可能有误；未完成数学TeX转写；禁止直接发给学生}
\begin{ocrdraft}
[教材p.122 / PDF第129页]
\par
比较下列各组中两个数的大小：
\par
（1)Iogz7.6和Iogz8.7，
\par
（2）
\par
Iogi7.6和log\ensuremath{\pm}8.7；
\par
（3）6和7伍>0且磊乒D；
\par
logo82和208
\par
（4）
\par
[教材p.123 / PDF第130页]
\par
解
\par
（1）因为函数y一log“在（0，十。。）上是增函数，且7．6<\&7，
\par
所以
\par
10g27，6<Iog28.7．
\par
（2）
\par
因为函数3-log丷在（0，十。。）上是减函数，且7，6<\&7，
\par
所以
\par
Iogi7.6>log上\&7．
\par
（3）
\par
当磊>1时，因为函数3-log辰在（0，十亠）上是增函数，且7，6<\&7，
\par
所以
\par
loga7.6<log\ocrglyph{201E}8.7；
\par
当0<<1时，因为函数在（0，十。。）上是减函数，
\par
所以
\par
log\ocrglyph{201E}7.6>logu8.7．
\par
（4）因为函数3-logo“在（0，十亠）上是减函数，
\par
所以
\par
logo.82<logo81一0．
\par
又因为20．8>0，所以logo.82<208
\end{ocrdraft}
\end{bankitem}
